\section{Computing comparison probabilities}\label{sec:oracle}

We now implement \texttt{Estimate} and prove Lemma~\ref{lem:estimation}.
The calculation retains a bounded neighborhood of the queried items and
computes the probability of the requested event exactly in that smaller
region. We first show why this probability approximates the one in the
full region, then give the finite calculation.

\subsection{The region used by the estimator}

Let $G$ be the in-node DAG and $K$ its feasible region. First compute the
effective endpoints $\ell_z,r_z$ using all of $G$, in temporary storage.
These endpoints preserve interval restrictions transmitted by paths that
may extend beyond the neighborhood we will retain.

For a threshold query $X_x<c$, retain the vertices at undirected distance
at most $\rho$ from $x$. For a pair query $X_x<X_y$, retain the union of
the two such neighborhoods. Keep every arc between retained vertices and
copy their effective endpoints from $G$. Write $V'$ for this set of
vertices, and let
\[
K'=\{z\in\mathbb R^{V'}:\ell_i<z_i<r_i\text{ for all }i\in V',\quad
                z_i<z_j\text{ for every retained arc }i\to j\}.
\]
The region $K'$ has positive volume: it contains the projection onto
$V'$ of every feasible assignment in $K$, including an open set.
For the requested event $E$, the estimator returns
\begin{equation}\label{eq:local-probability}
\widehat p=\frac{\vol(K'\cap E)}{\vol(K')}.
\end{equation}
If a directed path in $G$ already decides a pair query, or the threshold
lies outside the queried effective interval, the answer is exactly zero
or one and no volume calculation is needed.

The following theorem supplies the radius. Its dependence only on the
stated parameters is what makes the exact calculation possible in
polynomial time after the parameter choice in Section~\ref{sec:runtime}.

\begin{theorem}[Local probability estimate]\label{thm:local}
Let $d,h\ge1$ be integers. Suppose $G$ has degree at most $d$, its
directed paths have at most $h$ vertices, and $K\subseteq(0,1)^{V(G)}$
has positive volume. For every rational
$R\ge1$ and $0<\varepsilon<1/4$, there is a computable integer
$\rho=\rho(d,h,R,\varepsilon)$ with the following properties.
\begin{enumerate}
\item If $W_x\ge1/R$, the estimate in \eqref{eq:local-probability} for
$X_x<c$ differs from its probability in $K$ by at most $\varepsilon$.
\item If at least one of $W_x,W_y$ is at least $1/R$, the same holds for
$X_x<X_y$.
\end{enumerate}
The retained region has at most
\begin{equation}\label{eq:local-size}
b=2(d+1)^{\rho+1}
\end{equation}
coordinates. All its interval bounds are the effective bounds computed
in the full graph.
\end{theorem}

\subsection{Why distant coordinates have little effect}

We put uniform assignments on $K$ and $K'$ in a common random experiment
and bound the probability that they disagree on the query. For the proof,
we update the assignments one coordinate at a time. We trace differences
between corresponding coordinates through these updates and show that
every possible sequence of dependencies is likely to reduce them.

\begin{proof}[Proof of Theorem~\ref{thm:local}]
\emph{Updating the two assignments.}
Choose a topological order of $G$ and use its
restriction to $V'$ for the local region. One pass through this order
will be called a \emph{sweep}.

For an assignment uniform on either region, conditioning on the other
coordinates leaves a uniform distribution for $X_z$ on the interval between
\[
L_z^X=\max\bigl(\{\ell_z\}\cup\{X_u:u\to z\}\bigr),\qquad
R_z^X=\min\bigl(\{r_z\}\cup\{X_u:z\to u\}\bigr).
\]
Thus replacing $X_z$ by $(1-U_z)L_z^X+U_zR_z^X$, with $U_z$ uniform on
$(0,1)$, preserves the uniform distribution on the feasible region. This applies to each update
in either graph.

Start $X$ uniformly on $K$ and $Y$ uniformly on $K'$. At every retained
vertex, use the same fresh uniform number for their two updates. Use
independent fresh numbers at different vertices and sweeps. Both
assignments remain uniform on their own regions throughout this joint
experiment.

\emph{The decrease after one sweep.}
For $z\in V'$, write $\Delta_z=|X_z-Y_z|$. Superscripts old and new
refer to the start of a sweep and to the value after a coordinate's
update. Whenever all neighbors of $z$ are retained, the two updates use
the same interval bounds and give
\begin{equation}\label{eq:local-coupling}
\Delta_z^{\rm new}\le
(1-U_z)\max_{u\to z}\Delta_u^{\rm new}
+U_z\max_{z\to u}\Delta_u^{\rm old},
\end{equation}
where $U_z$ is their shared random number and an empty maximum is zero.
The predecessors have already been updated, while the successors retain
their old values.

Let $B_h(z)$ denote the vertices within undirected distance $h$ of $z$,
and suppose this whole set is retained. Follow a predecessor attaining
the first maximum in \eqref{eq:local-coupling}, and continue backwards
until reaching a source. The resulting path has at most $h$ vertices.
All its vertices and their successors lie in $B_h(z)$, so every inequality
used in this expansion is valid in both graphs. At the source, the
lower-endpoint term in \eqref{eq:local-coupling} is zero. It follows that
\begin{equation}\label{eq:local-one-sweep}
\Delta_z^{\rm new}\le(1-A_z^h)
                 \max_{u\in B_h(z)}\Delta_u^{\rm old},
\end{equation}
where $A_z$ is the minimum of $1-U_u$ over $z$ and all its ancestors.
Indeed, expansion along the selected path gives a factor
$1-\prod_u(1-U_u)$, which is at most $1-A_z^h$.

\emph{Following the possible dependencies.}
Both $B_h(z)$ and the set of ancestors used to define $A_z$ have at most
$M$ vertices, where
\begin{equation}\label{eq:local-moment}
M=(d+1)^{h+1}.
\end{equation}
Write $\Delta_z(s)$ for the discrepancy after $s$ sweeps and
$A_z(s)$ for the minimum used in sweep $s$.
Fix an even integer $t$. Repeatedly applying
\eqref{eq:local-one-sweep} traces a sequence
$z_t=z,z_{t-1},\ldots,z_0$, where $z_{s-1}\in B_h(z_s)$ at each step.
There are at most $M^t$ such sequences.

Suppose $B_{ht}(z)$ is retained. For $1\le s\le t$, every possible $z_s$
is at distance at most $(t-s)h$ from $z$, so $B_h(z_s)$ is retained too.
Thus every application of \eqref{eq:local-one-sweep} is valid in both
graphs. Since each initial discrepancy is at most one, we obtain
\begin{equation}\label{eq:local-contraction}
\Delta_z(t)\le
\max_{z_t=z,\;z_{s-1}\in B_h(z_s)}
       \prod_{s=1}^t\bigl(1-A_{z_s}(s)^h\bigr).
\end{equation}
At the outer boundary only the initial discrepancies occur in this
expansion.

We bound the probability that any of these products is large. Set
\[
a_0=\frac1{16M^3},\qquad \beta=1-a_0^h.
\]
Call a pair $(u,s)$ \emph{bad} if $A_u(s)<a_0$. A union bound over the
at most $M$ random numbers defining $A_u(s)$ gives
\[
\Prb\bigl((u,s)\text{ is bad}\bigr)\le Ma_0=\frac1{16M^2}.
\]
For every other pair, the corresponding factor in
\eqref{eq:local-contraction} is at most $\beta$.

For a fixed sequence $z_t,\ldots,z_0$, the events that its pairs are bad
involve independent random numbers from different sweeps. The probability that
at least $t/2$ of its pairs are bad is therefore at most
\[
2^t\left(\frac1{16M^2}\right)^{t/2}
=\left(\frac1{2M}\right)^t.
\]
Taking a union bound over the at most $M^t$ sequences gives probability
at most $2^{-t}$ that any sequence has this many bad pairs. Otherwise
every product in \eqref{eq:local-contraction} has at least $t/2$ factors
at most $\beta$. Consequently,
\begin{equation}\label{eq:local-interior-error}
\Prb\bigl(\Delta_z(t)>\beta^{t/2}\bigr)\le2^{-t}.
\end{equation}

\emph{From coordinate errors to comparison errors.}
Put $C=M+1$ and choose
\begin{align}
\delta&=\frac{\varepsilon}{8CR},&
\lambda&=\ceil{\log(1/\delta)},&
\mu&=\ceil{\log(4/\varepsilon)},\nonumber\\
t&=2\bigl(a_0^{-h}\lambda+\mu\bigr),&
\rho&=ht.\label{eq:local-radius}
\end{align}
Here $a_0^{-h}=(16M^3)^h$ is an integer, so $t$ is even. The inequality
$(1-v)^{-r}\ge1+rv$ gives $\beta^{a_0^{-h}}\le1/2$. Hence
$\beta^{t/2}\le2^{-\lambda}\le\delta$ and
$2^{-t}\le2^{-\mu}\le\varepsilon/4$. For each queried coordinate,
\eqref{eq:local-interior-error} now gives
\[
\Prb(\Delta_z(t)>\delta)\le\varepsilon/4.
\]
For an incomparable pair, choose the queried coordinate $x$ with
$W_x\ge1/R$. There are at most $M$ other coordinates comparable with
$x$. Corollary~\ref{cor:gap} gives
\[
\Prb(|X_x-X_y|\le2\delta)\le4CR\delta=\varepsilon/2.
\]
If both queried discrepancies are at most $\delta$, the two assignments
can disagree on $X_x<X_y$ only when $|X_x-X_y|\le2\delta$.
The disagreement probability is at most
$2(\varepsilon/4)+\varepsilon/2=\varepsilon$.
For a threshold, disagreement requires either $\Delta_x>\delta$ or
$|X_x-c|\le\delta$, so its probability is at most
$\varepsilon/4+2CR\delta=\varepsilon/2$.

The difference between the probabilities of an event is at most the
probability that its two indicators disagree. Since $X$ and $Y$ remain
uniform on $K$ and $K'$, these are precisely the probabilities being
compared. Finally, a radius-$\rho$ neighborhood has at most
$1+d+\cdots+d^\rho\le(d+1)^{\rho+1}$ vertices, proving
\eqref{eq:local-size}. All constants above are computable from the stated
parameters.
\end{proof}

\subsection{Calculating the local volumes exactly}\label{sec:exact-local-volume}

The number of retained coordinates is bounded independently of $n$ for
fixed adversary parameters. We can therefore enumerate their possible
orders and positions relative to the interval endpoints. This gives a
particularly direct exact calculation of \eqref{eq:local-probability}.

Suppose there are $r$ retained coordinates. Sort their interval endpoints,
together with $0,1$ and the threshold $c$ if present. Consecutive distinct
endpoints divide $(0,1)$ into at most $2r+2$ open segments. Enumerate a
segment for each coordinate and a permutation of all $r$ coordinates.
Keep a case exactly when the assigned segments lie inside the coordinates'
allowed intervals, their segment indices are nondecreasing along the
permutation, and that permutation respects all retained arcs. Each feasible assignment
outside a set of volume zero occurs in exactly one such case.

If segment $j$ has length $\Delta_j$ and contains $r_j$ coordinates in
the specified order, the case has volume
\begin{equation}\label{eq:exact-local-volume}
\prod_j\frac{\Delta_j^{r_j}}{r_j!}.
\end{equation}
Sum these volumes over all kept cases to obtain $\vol(K')$, and over
those satisfying the query to obtain $\vol(K'\cap E)$. A pair query is
decided by the permutation; a threshold query is decided by the assigned
segment. The denominator is positive, as established above.

\begin{lemma}[Cost of the local calculation]\label{lem:exact-local-volume}
For $r\le b$ coordinates numbered $1,\ldots,r$, and dyadic endpoints and threshold with at most
$\Lambda$ fractional binary digits, the two volumes and their ratio can
be computed in at most $C_0(b)(\Lambda+1)^2$ bit operations, for a
computable function $C_0$. The exponent of $\Lambda+1$ is independent of $b$.
\end{lemma}
\begin{proof}
There are at most $(2r+2)^r r!$ cases. Write each segment length as
$\Delta_j=a_j/2^\Lambda$. Every case volume in
\eqref{eq:exact-local-volume} has the common denominator
$2^{r\Lambda}r!$, because $\prod_j r_j!$ divides $r!$.
Its integer numerator is
\[
\frac{r!}{\prod_j r_j!}\prod_j a_j^{r_j}.
\]
The cases have disjoint interiors and lie in the unit cube. Both summed
numerators are therefore at most $2^{r\Lambda}r!$ and have
$O(r\Lambda+r\log(r+1))$ bits. Their ratio is the requested probability.
Sorting the endpoints, checking the cases, and forming their numerators
uses a number of operations depending only on $b$. Ordinary integer
arithmetic gives the stated bound.
\end{proof}

\begin{proof}[Proof of Lemma~\ref{lem:estimation}]
The strict inequalities defining the feasible region make it open, so
the existence of a feasible assignment gives positive volume.
Computing the effective endpoints in temporary storage preserves the
full feasible region by Lemma~\ref{lem:interior}. Apply
Theorem~\ref{thm:local} with $d=k$, $h=2H$, and
$\varepsilon=\eta$. The exact local calculation above returns the
probability to which that theorem applies, so it has the required
accuracy. It uses a finite deterministic enumeration with exact rational
arithmetic.
\end{proof}

The next section chooses the parameters as a function of $n$ and bounds
the running time over the entire interaction.
